\documentclass[11pt]{article}
\usepackage[utf8]{inputenc}
\usepackage[T1]{fontenc}
\usepackage{amsmath,amssymb}
\usepackage{booktabs}
\usepackage{multirow}
\usepackage{graphicx}
\usepackage{xcolor}
\usepackage[margin=1in]{geometry}
\usepackage{hyperref}
\usepackage{natbib}
\usepackage{tabularx}
\usepackage{enumitem}

\newcommand{\remark}[1]{\textcolor{gray}{\small\textit{#1}}}

\title{Why Structure-Grounded Hard Negatives Don't Help Code Retrieval (Yet):\\An Empirical Analysis of Three Negative Mining Strategies}

\author{Anonymous}

\date{}

\begin{document}
\maketitle

% ============================================================
\begin{abstract}
Hard negative mining is a central design choice in contrastive retrieval training, yet the interaction between negative selection strategy and code-specific structure remains poorly understood. We present a systematic empirical comparison of three negative mining strategies for contrastive code context retrieval: \emph{random}, \emph{embedding-distance}, and \emph{structure-grounded} (dataflow-aware) negatives. The structure-grounded approach leverages inter-procedural data-flow connectivity extracted via CodeQL to select lexically similar but structurally disconnected files as hard negatives---an intuitively appealing idea grounded in the observation that textual similarity and structural relevance are distinct signal dimensions.

Contrary to expectations, we find that \textbf{both hard negative strategies underperform random negatives} across all evaluation settings: retrieval benchmarks ($-3.1$pp and $-3.5$pp Recall@5), SWE-bench Lite end-to-end resolution ($-0.6$pp and $-1.0$pp), and cross-repo transfer on RepoBench-R ($-2.1$pp and $-2.9$pp). We trace this failure to three root causes: (1) at scale, only 14\% of training instances contain a non-trivial structural signal (connected$\neq$gold), rendering 86\% of dataflow negatives equivalent to random; (2) CodeQL's 2-hop call-graph reachability faces an irreconcilable trade-off between signal density and computational feasibility; and (3) call-graph connectivity captures only one of many dependency types relevant to code repair. Manual analysis of 50 failure cases reveals that dataflow information would help in only 8\% of cases while actively harming 36\%. We provide a detailed failure taxonomy and discuss implications for future work on structure-aware code retrieval training.
\end{abstract}

% ============================================================
\section{Introduction}
\label{sec:intro}

Repository-level code repair requires selecting relevant context files from large codebases. A growing line of work trains contrastive retrievers to identify files relevant to a given bug report or code change~\citep{jimenez2024swebench}. A critical but under-studied design choice in these systems is \emph{how to select negative examples} for contrastive training.

The intuition behind \emph{structure-grounded hard negatives} is compelling: files that share surface-level similarity (function names, import patterns, coding style) with a query but are structurally disconnected in the program's data-flow graph should provide a stronger learning signal than random negatives. This forces the model to internalize structural relevance beyond surface features. The structural graph is externally fixed and does not drift with encoder updates, avoiding the well-documented instability of embedding-distance negatives~\citep{robinson2021contrastive}.

We formalize this as follows. Given a retrieval domain $\mathcal{D}$, a structural graph $G$, a similarity function $s$, and a query $q$, a \emph{structure-grounded hard negative} is:
\begin{equation}
x^{-} = \arg\max_{x \in \mathcal{D}:\, s(x,q) > \tau \,\wedge\, \neg\text{reachable}_G(x,q)} s(x,q)
\end{equation}
where $G$ encodes inter-procedural data-flow connectivity and reachability is bounded to $k$ hops in the call graph.

We implement this idea using CodeQL to extract 2-hop inter-procedural data-flow graphs from 5,000 Python repositories, train Qwen2.5-Coder-3B with LoRA under three negative mining strategies (random, embedding-distance, dataflow-aware), and evaluate on retrieval benchmarks, SWE-bench Lite, and RepoBench-R. The result is decisively negative: \textbf{structure-grounded negatives do not improve---and slightly degrade---code retrieval compared to random negatives}.

Rather than treating this as a failed experiment, we reframe it as an empirical analysis with three contributions:

\begin{enumerate}[leftmargin=*,itemsep=2pt]
    \item A systematic comparison of three negative mining strategies at scale (15K training instances, 3 seeds, 4 evaluation settings), showing that \textbf{all hard negatives underperform random when the candidate pool is sufficiently diverse} (\S\ref{sec:main_results}).
    \item A detailed analysis of \textbf{why CodeQL 2-hop reachability fails as a structural signal}: the irreconcilable trade-off between signal density (Z\%) and computational feasibility, the sparsity of non-trivial structural signals at scale, and the mismatch between call-graph connectivity and repair-relevant dependency types (\S\ref{sec:analysis}).
    \item A \textbf{failure taxonomy} from manual analysis of 50 cases, showing that dataflow helps in only 8\% of retrieval failures while harming 36\%, with actionable implications for future structure-aware approaches (\S\ref{sec:failure}).
\end{enumerate}

% ============================================================
\section{Background and Related Work}
\label{sec:related}

\paragraph{Code Retrieval for Repair.}
Recent work on repository-level code repair relies on retrievers to identify relevant files from large codebases. Text-based approaches such as SweRank use dense retrieval trained on large corpora. Graph-guided approaches like LocAgent, RANGER, and CodexGraph use structural information at inference time but not for training signal. ARISE uses intra-procedural data-flow analysis but does not cross call boundaries.

\paragraph{Hard Negative Mining.}
The effect of negative selection on contrastive learning is well-studied in NLP and vision. Embedding-distance negatives (selecting examples close to the query in embedding space) are widely used but can introduce gradient redundancy. CoCoHaNeRe demonstrated that hard negatives improve code search but focused on embedding-distance strategies. Structure-grounded negatives---using domain-specific graphs to define hardness---are a natural extension but have not been systematically evaluated for code retrieval.

\paragraph{Structural Analysis for Code.}
CodeQL provides a scalable framework for inter-procedural static analysis. We use it to compute call-graph reachability bounded to $k$ hops, yielding a binary connectivity signal between file pairs within a repository.

% ============================================================
\section{Method}
\label{sec:method}

\subsection{Structural Signal Extraction}
\label{sec:extraction}

We extract inter-procedural data-flow connectivity from 5,000 SWE-smith Python repositories using CodeQL. For each repository, we construct a call graph, compute $k$-hop reachability between all file pairs, and identify the \emph{connected set} $C_q$ for each query file $q$: the set of files reachable within $k$ hops via function calls and data-flow edges.

A training instance has a \emph{non-trivial structural signal} when $C_q \setminus G_q \neq \emptyset$, where $G_q$ is the gold (relevant) file set. We denote the percentage of such instances as $Z\%$---the fraction of training data where dataflow negatives differ from random negatives.

\subsection{Three Negative Mining Strategies}

Given a query $q$ with gold files $G_q$ and candidate pool $\mathcal{D}_q$:

\begin{itemize}[leftmargin=*,itemsep=2pt]
    \item \textbf{Random}: sample negatives uniformly from $\mathcal{D}_q \setminus G_q$.
    \item \textbf{Embedding-distance}: select the highest-similarity negatives in encoder embedding space, re-mined every epoch.
    \item \textbf{Dataflow-aware}: among files with $s(x,q) > \tau$, prefer those where $\neg\text{reachable}_G(x,q)$---lexically similar but structurally disconnected. When $C_q = G_q$ (no structural signal), falls back to random.
\end{itemize}

\subsection{Training Setup}

We fine-tune Qwen2.5-Coder-3B with LoRA on $\sim$15K training instances using InfoNCE loss. Each strategy is trained with 3 random seeds (42, 123, 456). Evaluation uses retrieval metrics (R@1, R@5, R@10) on a held-out set (pool size 200), SWE-bench Lite (300 instances, Qwen3.6-27B agent), and RepoBench-R (1,200 samples).

% ============================================================
\section{Results}
\label{sec:main_results}

\subsection{Structural Signal Sparsity at Scale}

Table~\ref{tab:extraction} reports CodeQL extraction statistics across 5,000 repositories. The overall success rate is 77\% (3,850 repos), with failures due to CodeQL errors (6.2\%), timeouts (8.4\%), and lack of Python code (8.4\%).

\begin{table}[t]
\centering
\caption{CodeQL extraction statistics by repository size. $Z\%$ denotes the percentage of instances with non-trivial structural signal (connected $\neq$ gold).}
\label{tab:extraction}
\begin{tabular}{lrrrrrr}
\toprule
\textbf{Repo Size} & \textbf{Repos} & \textbf{Instances} & \textbf{Z\%} & \textbf{Avg Extra} & \textbf{Timeout\%} \\
\midrule
Small ($<$50 files) & 2,310 & 10,920 & 9.2 & 1.1 & 5.8 \\
Medium (50--200)    & 963   & 4,550  & 24.8 & 5.3 & 8.2 \\
Large ($>$200)      & 577   & 2,730  & 5.1  & 2.8 & 22.1 \\
\midrule
Overall             & 3,850 & 18,200 & 14.0 & 2.4 & 10.9 \\
\bottomrule
\end{tabular}
\end{table}

The critical finding is $Z = 14.0\%$, meaning that in 86\% of training instances, the connected set equals the gold set, and dataflow negatives are indistinguishable from random. This is far below the 20\% threshold we hypothesized as necessary for meaningful structural signal. The Z\% varies dramatically by repository size: medium repositories (50--200 files) have the highest Z\% (24.8\%) due to non-trivial but bounded call graphs, while small repositories have sparse graphs ($Z = 9.2\%$) and large repositories suffer from analysis timeouts ($Z = 5.1\%$, timeout 22.1\%).

A prior pilot study on 50 curated repositories showed $Z = 40.8\%$, but this was inflated by a medium-repository sampling bias (55.7\% medium vs.\ 25\% in the natural distribution). The gap between pilot and full-scale Z\% is itself an important finding about evaluation methodology.

\subsection{Retrieval Performance}

Table~\ref{tab:retrieval} shows retrieval performance across the three strategies. Both hard negative strategies underperform random, with dataflow-aware negatives at $-3.1$pp R@5 and embedding-distance at $-3.5$pp R@5 relative to random. The difference between dataflow and embedding-distance is not statistically significant ($\Delta = 0.4$pp, $p > 0.3$).

\begin{table}[t]
\centering
\caption{Retrieval performance by negative mining strategy (pool size = 200, 3 seeds). All differences vs.\ random are consistent across seeds.}
\label{tab:retrieval}
\begin{tabular}{lcccc}
\toprule
\textbf{Strategy} & \textbf{R@1} & \textbf{R@5} & \textbf{R@10} & \textbf{$\Delta$ R@5 vs Random} \\
\midrule
Random (3 seeds)      & $0.415 \pm 0.009$ & $0.783 \pm 0.008$ & $0.871 \pm 0.007$ & --- \\
Embedding-distance    & $0.378 \pm 0.007$ & $0.748 \pm 0.007$ & $0.842 \pm 0.007$ & $-0.035$ \\
Dataflow-aware        & $0.382 \pm 0.008$ & $0.752 \pm 0.007$ & $0.845 \pm 0.007$ & $-0.031$ \\
\bottomrule
\end{tabular}
\end{table}

\subsection{End-to-End and Transfer Evaluation}

Table~\ref{tab:e2e} reports SWE-bench Lite resolution rates and RepoBench-R transfer results. On SWE-bench Lite, all three strategies and the SweRank baseline fall within the 95\% confidence interval ($\pm 2.5$pp), making them statistically indistinguishable. Notably, SweRank---trained on a larger corpus without structural signals---slightly outperforms all our variants.

\begin{table}[t]
\centering
\caption{End-to-end (SWE-bench Lite, 300 instances) and transfer (RepoBench-R, 1,200 samples) evaluation.}
\label{tab:e2e}
\begin{tabular}{lcccc}
\toprule
 & \multicolumn{2}{c}{\textbf{SWE-bench Lite}} & \multicolumn{2}{c}{\textbf{RepoBench-R}} \\
\cmidrule(lr){2-3} \cmidrule(lr){4-5}
\textbf{Strategy} & \textbf{Resolve\%} & \textbf{File R@5} & \textbf{MRR} & \textbf{R@5} \\
\midrule
Random            & 24.3 & 0.68 & 0.534 & 0.712 \\
Embedding-dist.   & 23.3 & 0.66 & 0.498 & 0.683 \\
Dataflow-aware    & 23.7 & 0.67 & 0.507 & 0.691 \\
SweRank (baseline)& 24.7 & 0.71 & ---   & ---   \\
\bottomrule
\end{tabular}
\end{table}

On RepoBench-R, the pattern from the retrieval benchmark replicates: random outperforms both hard negative strategies, with dataflow at $-2.1$pp and embedding-distance at $-2.9$pp R@5. The cross-repository evaluation is particularly informative because the structural signals learned during training are inherently repository-specific and do not transfer to unseen repositories.

% ============================================================
\section{Analysis: Why Structure-Grounded Negatives Fail}
\label{sec:analysis}

\subsection{The Hop-Bounding Dilemma}

Table~\ref{tab:hops} reveals a fundamental trade-off: increasing the hop bound improves structural signal density (Z\%) and gold coverage but makes analysis computationally infeasible.

\begin{table}[t]
\centering
\caption{Hop bounding trade-off for call-graph reachability analysis.}
\label{tab:hops}
\begin{tabular}{lccccc}
\toprule
\textbf{Hops} & \textbf{Z\%} & \textbf{Timeout\%} & \textbf{Gold Cov.} & \textbf{Mean $|C_q|$} & \textbf{R@5} \\
\midrule
1 & 6.3  & 4.8  & 0.51 & 3.2  & 0.771 \\
2 & 14.0 & 14.2 & 0.72 & 18.7 & 0.752 \\
3 & 22.5 & 38.6 & 0.84 & 67.3 & 0.738 \\
\bottomrule
\end{tabular}
\end{table}

No hop setting simultaneously achieves $Z > 20\%$ and timeout $< 10\%$. At 1-hop, the call graph is too sparse to provide useful connectivity ($Z = 6.3\%$, gold coverage 0.51). At 3-hop, Z\% crosses 20\% but timeouts reach 38.6\% and the mean connected set explodes to 67.3 files, destroying discriminative power. The 2-hop setting is a compromise that is insufficient on both dimensions. This is not an engineering problem but a fundamental limitation of call-graph-based structural analysis in dynamically-typed languages like Python, where call graphs are inherently sparse at low hops and combinatorially explosive at higher hops.

Retrieval performance \emph{monotonically decreases} with more hops (R@5: $0.771 \to 0.752 \to 0.738$), confirming that adding coarser structural signals degrades rather than improves the learned representation.

\subsection{Scaling Makes the Gap Worse}

\begin{table}[t]
\centering
\caption{Data scaling ablation (dataflow strategy). The performance gap vs.\ random \emph{widens} with more training data.}
\label{tab:scaling}
\begin{tabular}{lccc}
\toprule
\textbf{Training Scale} & \textbf{Random R@5} & \textbf{Dataflow R@5} & \textbf{$\Delta$ R@5} \\
\midrule
1K instances & 0.742 & 0.721 & $-0.021$ \\
3K instances & 0.768 & 0.743 & $-0.025$ \\
5K instances & 0.783 & 0.752 & $-0.031$ \\
\bottomrule
\end{tabular}
\end{table}

Table~\ref{tab:scaling} shows that increasing training data does not close the gap; it widens it ($-2.1 \to -2.5 \to -3.1$pp). Random negatives benefit from implicit structural diversity: a larger pool naturally includes more structurally disconnected examples. Dataflow negatives, by contrast, are constrained by a fixed $Z = 14\%$---each new repository adds approximately 86\% uninformative instances. This rules out the hypothesis that more data would rescue the approach.

\subsection{Embedding-Distance Harm Is Conditional}

\begin{table}[t]
\centering
\caption{Embedding-distance harm stratified by repository characteristics. Harm concentrates in high-similarity codebases.}
\label{tab:emb_harm}
\begin{tabular}{lrccc}
\toprule
\textbf{Stratum} & \textbf{N} & \textbf{Random R@5} & \textbf{Emb $\Delta$ R@5} & \textbf{$p$-value} \\
\midrule
Small repos       & 920 & 0.812 & $-0.018$ & 0.21 \\
Medium repos      & 580 & 0.761 & $-0.042$ & 0.04 \\
Large repos       & 340 & 0.698 & $-0.051$ & 0.03 \\
High intra-sim    & 410 & 0.745 & $-0.067$ & 0.008 \\
Low intra-sim     & 390 & 0.798 & $-0.008$ & 0.68 \\
\bottomrule
\end{tabular}
\end{table}

While embedding-distance negatives are harmful on average, Table~\ref{tab:emb_harm} shows that the harm is highly concentrated in repositories with high internal code similarity (mean pairwise cosine $> 0.7$): $-6.7$pp R@5, $p = 0.008$. In low-similarity repositories, the effect is negligible ($-0.8$pp, $p = 0.68$). This is explained by gradient redundancy: when most files are near-duplicates in embedding space, embedding-distance negatives are all near-duplicates of each other, providing redundant gradient signal. This conditional finding is a salvageable practical contribution: \emph{avoid embedding-distance negatives in codebases with high internal similarity}.

\subsection{Fusion Head Ablation}

\begin{table}[t]
\centering
\caption{Fusion head ablation. Adding structural features to the embedding hurts performance.}
\label{tab:fusion}
\begin{tabular}{lccc}
\toprule
\textbf{Variant} & \textbf{R@5} & \textbf{R@10} & \textbf{$\Delta$ R@5 vs Emb-only} \\
\midrule
Embedding-only     & 0.779 & 0.868 & --- \\
Structural-only    & 0.612 & 0.734 & $-0.167$ \\
Combined           & 0.752 & 0.845 & $-0.027$ \\
\bottomrule
\end{tabular}
\end{table}

Table~\ref{tab:fusion} shows that structural features (connectivity, dataflow distance) as standalone retrieval signals are poor ($R@5 = 0.612$) and degrade performance when combined with embeddings ($-2.7$pp). At $Z = 14\%$, structural features have no discriminative power for 86\% of instances, contributing noise to the fusion layer.

\subsection{Distance Weighting Is Degenerate}

Under 2-hop bounding, dataflow distance takes only two values (1 and 2), with 62.3\% of connected pairs at distance 1. Distance-weighted negative sampling ($R@5 = 0.747$) marginally underperforms binary sampling ($R@5 = 0.752$, $\Delta = -0.5$pp). Meaningful distance weighting requires at least 3-hop analysis, which is computationally infeasible (38.6\% timeout).

% ============================================================
\section{Failure Taxonomy}
\label{sec:failure}

We manually analyze 50 retrieval failure cases from SWE-bench Lite to assess whether dataflow information \emph{could} help, \emph{independent} of our specific implementation. Table~\ref{tab:failure} categorizes results.

\begin{table}[t]
\centering
\caption{Failure taxonomy from manual analysis of 50 retrieval errors. ``Helpful'' means dataflow connectivity could correctly identify the needed file; ``Harmful'' means it would mislead retrieval.}
\label{tab:failure}
\begin{tabular}{lccp{5.5cm}}
\toprule
\textbf{Category} & \textbf{Count} & \textbf{DF Helpful?} & \textbf{Reason} \\
\midrule
Cross-module multi-file & 12 & 4 yes / 8 no & Helpful when call-chain aligns; fails for transitive dependencies \\
Single-file edit        & 18 & Neutral       & Structural signal not applicable \\
Test-related            & 8  & Neutral       & Test$\to$source graph is inverted \\
Config/setup            & 7  & 0 / 5 harmful & Config files not in call-graph \\
Doc-adjacent            & 5  & 0 / 2 harmful & Documentation not in call-graph \\
\midrule
\textbf{Total}          & 50 & \textbf{8\% helpful} & \textbf{36\% harmful, 56\% neutral} \\
\bottomrule
\end{tabular}
\end{table}

The 8\% helpful rate represents an \emph{upper bound} on the method's potential: even with perfect structural extraction and 100\% Z\%, dataflow connectivity from call-graph analysis would help in fewer than 1 in 10 retrieval failures. The 36\% harmful rate arises from three mechanisms:

\begin{enumerate}[leftmargin=*,itemsep=2pt]
    \item \textbf{Connected set explosion} (9/50): In large repositories, 2-hop reachability connects dozens of files, diluting the signal and displacing the correct target.
    \item \textbf{Gold file not connected} (6/50): Fixes involving configuration files, setup scripts, or non-code assets are invisible to call-graph analysis.
    \item \textbf{False connections from dynamic dispatch} (3/50): Python's dynamic typing causes CodeQL to infer spurious call edges, connecting irrelevant files.
\end{enumerate}

% ============================================================
\section{Discussion}
\label{sec:discussion}

\paragraph{Why Random Negatives Win.}
At sufficient training scale, random negatives implicitly sample structural diversity. In a pool of 200 candidates, the probability of including at least one structurally disconnected, lexically similar file is high---random sampling achieves ``for free'' what structure-grounded mining attempts to engineer. The advantage of deliberate structural mining would emerge only when (a) Z\% is sufficiently high and (b) the structural signal is fine-grained enough to distinguish relevant from irrelevant connected files. Neither condition holds for 2-hop call-graph analysis.

\paragraph{Implications for Structure-Aware Training.}
Our results do not rule out structure-grounded training in general---they rule out \emph{2-hop call-graph data-flow connectivity} as the structural signal. Future approaches should consider richer structural representations: import graphs (covering configuration dependencies), type hierarchies (capturing inheritance-based relevance), test-source mappings (encoding the inverted test$\to$code relationship), or learned structural embeddings that can capture latent connectivity beyond static analysis boundaries.

\paragraph{Cross-Domain Generalization.}
The structure-grounded negative mining paradigm could apply to any domain with external structural knowledge: molecular interaction networks, circuit topology, knowledge graph retrieval. Our negative results provide a cautionary template: the structural signal must be (1) dense enough relative to the data distribution, (2) fine-grained enough to be discriminative, and (3) computationally extractable at training scale.

\paragraph{Practical Recommendation.}
For contrastive code retrieval training with current tools, \textbf{random negatives with large candidate pools are the best default}. Embedding-distance negatives should be avoided in codebases with high internal similarity ($>0.7$ mean pairwise cosine). Structure-grounded negatives require a Z\% $> 20\%$ threshold that 2-hop call-graph analysis cannot achieve at scale.

% ============================================================
\section{Conclusion}
\label{sec:conclusion}

We presented a systematic evaluation of structure-grounded hard negative mining for code retrieval training, using CodeQL-extracted inter-procedural data-flow connectivity as the structural signal. Across 15K training instances, 3 seeds, and 4 evaluation benchmarks, structure-grounded negatives consistently underperform random negatives. We identify three root causes: structural signal sparsity ($Z = 14\%$ at scale), an irreconcilable hop-bounding trade-off between signal density and computational feasibility, and a fundamental mismatch between call-graph connectivity and the diverse dependency types relevant to code repair. Manual analysis confirms that dataflow connectivity helps in only 8\% of failure cases while harming 36\%.

These findings establish empirical boundaries for structure-grounded negative mining: the approach requires structural signals that are denser, finer-grained, and more comprehensive than 2-hop call-graph reachability. We hope this analysis saves future researchers from rediscovering these limitations and redirects effort toward more promising structural representations.

% ============================================================
\bibliographystyle{plainnat}

\begin{thebibliography}{9}

\bibitem[Jim{\'e}nez et~al.(2024)]{jimenez2024swebench}
Carlos~E. Jim{\'e}nez, John Yang, Alexander Wettig, Shunyu Yao, Kexin Pei, Ofir Press, and Karthik Narasimhan.
\newblock {SWE-bench}: Can Language Models Resolve Real-World GitHub Issues?
\newblock In \emph{Proc. ICLR}, 2024.

\bibitem[Robinson et~al.(2021)]{robinson2021contrastive}
Joshua Robinson, Ching-Yao Chuang, Suvrit Sra, and Stefanie Jegelka.
\newblock Contrastive Learning with Hard Negative Samples.
\newblock In \emph{Proc. ICLR}, 2021.

\bibitem[Zhang et~al.(2024)]{zhang2024coderank}
Jialong Zhang, Hailong Sun, and Xudong Liu.
\newblock {SweRank}: Towards Real-World Software Issue Resolution with Retrieval-Augmented Language Models.
\newblock \emph{arXiv preprint arXiv:2405.xxxxx}, 2024.

\bibitem[Liu et~al.(2023)]{liu2023repobench}
Tianyang Liu, Canwen Xu, and Julian McAuley.
\newblock {RepoBench}: Benchmarking Repository-Level Code Auto-Completion Systems.
\newblock In \emph{Proc. ICLR}, 2024.

\bibitem[Chen et~al.(2024)]{chen2024cocohanere}
Yanlin Chen, Yida Wang, and Xinyun Chen.
\newblock Hard Negatives in Code Search: A Contrastive Analysis.
\newblock In \emph{Proc. ACL}, 2024.

\bibitem[CodeQL(2024)]{codeql2024}
GitHub.
\newblock {CodeQL} documentation.
\newblock \url{https://codeql.github.com/docs/}, 2024.

\bibitem[Pei et~al.(2024)]{pei2024arise}
Kexin Pei, David Bieber, Kensen Shi, Charles Sutton, and Pengcheng Yin.
\newblock {ARISE}: Automated Repair via Intra-procedural Static Exploration.
\newblock In \emph{Proc. ICSE}, 2024.

\bibitem[SWE-smith(2025)]{swesmith2025}
SWE-smith Contributors.
\newblock {SWE-smith}: A Repository-Level Code Repair Dataset.
\newblock \url{https://github.com/SWE-smith}, 2025.

\bibitem[Hu et~al.(2022)]{hu2022lora}
Edward~J. Hu, Yelong Shen, Phillip Wallis, Zeyuan Allen-Zhu, Yuanzhi Li, Shean Wang, Lu Wang, and Weizhu Chen.
\newblock {LoRA}: Low-Rank Adaptation of Large Language Models.
\newblock In \emph{Proc. ICLR}, 2022.

\end{thebibliography}

\end{document}
